% ---------------------------------------------------------------------
%  Notation
%
%  Suggested placement: in the frontmatter, after the preface and before
%  the table of contents. In thesis.tex, add (inside \frontmatter):
%
%      % ---------------------------------------------------------------------
%  Notation
%
%  Suggested placement: in the frontmatter, after the preface and before
%  the table of contents. In thesis.tex, add (inside \frontmatter):
%
%      % ---------------------------------------------------------------------
%  Notation
%
%  Suggested placement: in the frontmatter, after the preface and before
%  the table of contents. In thesis.tex, add (inside \frontmatter):
%
%      % ---------------------------------------------------------------------
%  Notation
%
%  Suggested placement: in the frontmatter, after the preface and before
%  the table of contents. In thesis.tex, add (inside \frontmatter):
%
%      \include{frontmatter/notation}
%
%  This uses only the longtable environment from the `longtable` package.
%  Add \usepackage{longtable} to the preamble if it is not already loaded
%  (amsmath/amssymb are already loaded in thesis.tex).
% ---------------------------------------------------------------------

\chapter*{Notation}
\addcontentsline{toc}{chapter}{Notation}
\markboth{Notation}{Notation}

The following list summarizes the notation used throughout this thesis. Symbols
are grouped by theme and, within each group, ordered roughly by first
appearance. Page or equation references are given where a symbol is formally
defined. Throughout, the circle angle is written $\theta$ in the introductory
linear-model discussion (Chapter~\ref{cha:intro}) and $\varphi$ from the
methodology onward (Chapters~\ref{cha:methodology}--\ref{cha:results}); the two
denote the same angular coordinate on $S^1$.

\vspace{1em}

\begin{longtable}{p{0.26\textwidth} p{0.68\textwidth}}

% ----------------------------------------------------------------
\multicolumn{2}{l}{\textbf{Sets, spaces, and measures}}\\[0.3em]
\endfirsthead
\multicolumn{2}{l}{\small\itshape (Notation, continued)}\\[0.3em]
\endhead

$\mathbb{R},\ \mathbb{Z}$ & The real numbers and the integers. \\[0.4em]
$\mathbb{R}^{d}$ & $d$-dimensional Euclidean space. \\[0.4em]
$S^1$ & The unit circle, identified with the angular domain $\Omega=[0,2\pi)$ and embedded in $\mathbb{R}^2$ via $x(\varphi)=(\cos\varphi,\sin\varphi)$. \\[0.4em]
$S^{d-1}$ & The unit sphere in $\mathbb{R}^d$ (the $(d-1)$-sphere). \\[0.4em]
$\mathcal{X},\ \mathcal{Y}$ & Input space ($\mathcal{X}\subseteq\mathbb{R}^d$) and target space (e.g.\ $\mathcal{Y}=\mathbb{R}^n$ in regression, $\mathcal{Y}=\{1,\dots,C\}$ in classification). \\[0.4em]
$\rho$ & Joint data distribution over $\mathcal{X}\times\mathcal{Y}$; also used for the uniform input distribution. \\[0.4em]
$\mu$ & Uniform probability measure on $S^1$, $d\mu(\varphi)=d\varphi/(2\pi)$. \\[0.4em]
$L^2(\mu)$ & Space of square-integrable functions on $S^1$ with respect to $\mu$, with inner product $\langle\cdot,\cdot\rangle_{L^2(\mu)}$. \\[0.4em]
$\mathcal{H}$ & A reproducing kernel Hilbert space (RKHS). \\[0.4em]
$\mathcal{H}_\theta$ & Tangent parameter space ($\mathbb{R}^p$ in finite dimension; its limit in the infinite-width case). \\[0.4em]
$\Omega$ & The angular domain $[0,2\pi)$; in the mean-field discussion, the neuron parameter space $\mathbb{R}\times\mathbb{R}^d$. \\[0.4em]
$\|\cdot\|_2,\ \|\cdot\|_F,\ \|\cdot\|_{\max}$ & Euclidean ($\ell^2$) norm, Frobenius norm, and entrywise max norm, respectively. \\[0.4em]
$\langle\cdot,\cdot\rangle$ & Euclidean / Hilbert-space inner product. \\[0.4em]
$\mathbb{E},\ \mathbb{P}$ & Expectation and probability. \\[0.4em]
$\mathcal{N}(0,\sigma^2)$ & Gaussian distribution with mean $0$ and variance $\sigma^2$; $I_d$ denotes the $d\times d$ identity. \\[0.6em]

% ----------------------------------------------------------------
\multicolumn{2}{l}{\textbf{Network, data, and training}}\\[0.3em]
$f(x;\theta),\ f_\theta$ & Neural network output at input $x$ with parameter vector $\theta$. \\[0.4em]
$f_m(x;\theta)$ & One-hidden-layer ReLU network of width $m$ in NTK parameterization, $f_m(x;\theta)=\tfrac{1}{\sqrt{m}}\sum_{i=1}^m a_i\,\sigma(w_i^\top x+b_i)$. \\[0.4em]
$\theta$ & Parameter vector, $\theta\in\mathbb{R}^p$; $\theta_0$ its value at initialization. \\[0.4em]
$p$ & Number of network parameters. \\[0.4em]
$m$ & Network width (number of hidden units); also $w$ as a width label in figures. \\[0.4em]
$d$ & Input dimension; also the dimension $d=2K+1$ of the truncated Fourier block $\mathcal{H}_K$. \\[0.4em]
$a_i,\ w_i,\ b_i$ & Output weight, input weight vector, and bias of hidden unit $i$; initialized $a_i\sim\mathcal{N}(0,1)$, $w_i\sim\mathcal{N}(0,I_2)$, $b_i(0)=0$. \\[0.4em]
$v_\alpha,\ w_\alpha$ & Second-layer and first-layer parameters of neuron $\alpha$ in the generic two-layer network. \\[0.4em]
$\sigma,\ \varphi(\cdot)$ & Activation function; $\sigma(t)=t_+=\max\{t,0\}$ is ReLU ($\varphi$ also denotes a generic activation, and elsewhere the circle angle). \\[0.4em]
$W,\ v$ & Hidden-layer weight matrix and output-weight vector of a two-layer network. \\[0.4em]
$\odot$ & Hadamard (elementwise) product. \\[0.4em]
$\mathcal{D}=\{(x_i,y_i)\}_{i=1}^n$ & Training dataset of $n$ i.i.d.\ samples. \\[0.4em]
$n$ & Number of training samples / sampled points. \\[0.4em]
$X$ & Data matrix ($X\in\mathbb{R}^{d\times n}$, columns $x_i$). \\[0.4em]
$y$ & Vector / function of targets. \\[0.4em]
$\ell(\cdot,\cdot)$ & Loss function. \\[0.4em]
$\mathcal{L},\ L$ & Empirical (least-squares) risk; $\mathcal{L}_\rho$ denotes the population risk. \\[0.4em]
$\eta$ & Gradient-descent step size / learning rate. \\[0.4em]
$t,\ k$ (subscript) & Training time / iteration index. \\[0.4em]
$\nabla_\theta$ & Gradient with respect to $\theta$; $\dot{(\cdot)}$ denotes time derivative under gradient flow. \\[0.4em]
$\theta^\star,\ w^\star,\ f^\star$ & Minimizers / fixed points of the relevant objective. \\[0.4em]
$A^{+},\ A^{\ast}$ & Moore--Penrose pseudoinverse and adjoint of an operator/matrix $A$. \\[0.4em]
$\mathrm{range}(\cdot),\ \ker(\cdot)$ & Range (column space) and kernel (null space) of an operator. \\[0.4em]
$\lambda_{\max}(\cdot)$ & Largest eigenvalue of a matrix. \\[0.6em]

% ----------------------------------------------------------------
\multicolumn{2}{l}{\textbf{Neural Tangent Kernel and operators}}\\[0.3em]
$\phi(x),\ \Phi_t(x)$ & Tangent feature map $\nabla_\theta f(x;\theta_0)$ (and its time-$t$ version). \\[0.4em]
$\Theta_0(x,x')$ & Neural Tangent Kernel at initialization, $\nabla_\theta f(x;\theta_0)^\top\nabla_\theta f(x';\theta_0)$. \\[0.4em]
$\Theta_t,\ \Theta_t^{(m)}$ & NTK at training time $t$ (the empirical, width-$m$ kernel). \\[0.4em]
$\Theta_0^{(m)}$ & Finite-width empirical NTK at initialization. \\[0.4em]
$\bar\Theta,\ \Theta$ & Limiting (infinite-width) deterministic NTK; also written as the Gram matrix $\bar\Theta=(\bar\Theta(x_i,x_j))_{i,j}$. \\[0.4em]
$\Theta(\delta)$ & Limiting NTK on $S^1$ as a function of the angular distance $\delta=d(\varphi,\psi)\in[0,\pi]$. \\[0.4em]
$\Theta^{(\mathrm{nb})}$ & Limiting NTK with the trainable hidden-bias contribution removed. \\[0.4em]
$\Theta^{(r)},\ \psi_r$ & One-neuron kernel contribution and one-neuron tangent feature of neuron $r$. \\[0.4em]
$J_0$ & Jacobian of network outputs w.r.t.\ parameters at initialization; $\bar\Theta=J_0 J_0^\ast$. \\[0.4em]
$T_{\bar\Theta},\ T,\ T_\infty$ & Integral operator on $L^2(\mu)$ induced by the limiting kernel; $T_\infty=\lim_{m\to\infty}T_0^{(m)}$. \\[0.4em]
$T_t^{(m)},\ T_0^{(m)}$ & Empirical tangent integral operator at time $t$ and at initialization (frozen operator). \\[0.4em]
$T^{(r)}$ & One-neuron integral operator; $T_0^{(m)}=\tfrac1m\sum_{r=1}^m T^{(r)}$. \\[0.4em]
$T_p$ & Density-weighted integral operator for a non-uniform input density $p(\theta)$. \\[0.4em]
$\kappa$ & The $2\pi$-periodic kernel profile with $\bar\Theta(\theta,\theta')=\kappa(\theta-\theta')$; also a uniform kernel bound and (in the conditioning discussion) a matrix condition number. \\[0.4em]
$r_t,\ r$ & Residual $f_t-y$ (function or vector). \\[0.6em]

% ----------------------------------------------------------------
\multicolumn{2}{l}{\textbf{Fourier basis, spectrum, and subspaces}}\\[0.3em]
$\theta,\ \varphi,\ \gamma$ & Angular coordinate on $S^1$; $\gamma$ is used as the angle variable in target functions. \\[0.4em]
$\phi_q(\theta)=e^{iq\theta}$ & Complex Fourier mode of frequency $q\in\mathbb{Z}$. \\[0.4em]
$\phi_0,\ \phi_{k,c},\ \phi_{k,s}$ & Real Fourier basis of $L^2(\mu)$: $\phi_0=1$, $\phi_{k,c}=\sqrt{2}\cos(k\varphi)$, $\phi_{k,s}=\sqrt{2}\sin(k\varphi)$. \\[0.4em]
$\psi_j$ & Generic orthonormal eigenfunction of an integral operator. \\[0.4em]
$k,\ q$ & Fourier frequency index; $k(p)$ is the frequency of basis function $\phi_p$. \\[0.4em]
$K$ & Frequency cutoff defining the retained low-frequency block. \\[0.4em]
$\mathcal{H}_K$ & Truncated Fourier subspace $\operatorname{span}\{\phi_0,\phi_{k,c},\phi_{k,s}:1\le k\le K\}$, of dimension $d=2K+1$. \\[0.4em]
$\mathcal{F}_k$ & Frequency-$k$ subspace: $\mathcal{F}_0=\operatorname{span}\{\phi_0\}$, $\mathcal{F}_k=\operatorname{span}\{\phi_{k,c},\phi_{k,s}\}$ for $k\ge1$. \\[0.4em]
$d_k$ & Dimension of $\mathcal{F}_k$: $d_0=1$, $d_k=2$ for $k\ge1$. \\[0.4em]
$\lambda_q,\ \lambda_k$ & Fourier eigenvalue of $T$ at frequency $q$ (resp.\ $k$); $\lambda_p:=\lambda_{k(p)}$. \\[0.4em]
$\lambda_k^{(\mathrm{nb})}$ & Fourier eigenvalue of the no-bias kernel $\Theta^{(\mathrm{nb})}$. \\[0.4em]
$\Lambda$ & Diagonal matrix of eigenvalues $\operatorname{diag}(\lambda_p)$; in the linear model, $A=Q\Lambda Q^\top$. \\[0.4em]
$\Lambda_K$ & Continuum reference block $P_K T_\infty P_K=\operatorname{diag}(\lambda_0,\lambda_1,\lambda_1,\dots,\lambda_K,\lambda_K)$. \\[0.4em]
$\alpha_p(t)$ & Amplitude of the residual along Fourier mode $\phi_p$; decays as $e^{-\lambda_p t}\alpha_p(0)$. \\[0.4em]
$a_j(t)$ & Coefficient of the residual along eigenfunction $\psi_j$ in the continuum operator analysis. \\[0.6em]

% ----------------------------------------------------------------
\multicolumn{2}{l}{\textbf{Finite-sample objects}}\\[0.3em]
$\varphi_1,\dots,\varphi_n$ & I.i.d.\ angles sampled uniformly on $[0,2\pi)$. \\[0.4em]
$\Phi$ & Sampled Fourier feature matrix, $\Phi\in\mathbb{R}^{n\times d}$, $\Phi_{i,p}=\phi_p(\varphi_i)$; $\Phi_p$ its $p$-th column. \\[0.4em]
$A$ & Normalized sampled kernel matrix, $A_{ij}=\tfrac1n\Theta(\varphi_i-\varphi_j)$. \\[0.4em]
$G$ & Sampled Fourier Gram matrix, $G=\tfrac1n\Phi^\top\Phi$. \\[0.4em]
$H$ & Compressed operator matrix, $H=\tfrac1n\Phi^\top A\Phi$; $\mathbb{E}[H]=\Lambda^{(n)}$. \\[0.4em]
$\lambda_p^{(n)}$ & Finite-sample corrected Fourier eigenvalue, $(1-\tfrac1n)\lambda_p+\tfrac{\Theta(0)}{n}$. \\[0.4em]
$\Lambda^{(n)}$ & Diagonal matrix $\operatorname{diag}(\lambda_p^{(n)})$ of corrected eigenvalues. \\[0.4em]
$C,\ C_n^{(K)}$ & Coefficient-level restricted operator $G^{-1}H$ on the retained block; compared against $\Lambda^{(n)}$. \\[0.4em]
$P_{\mathrm{th}},\ P_{\mathrm{emp}}$ & Theory and empirical Fourier-block preconditioners. \\[0.4em]
$B_{\mathrm{base}},\ B_{\mathrm{th}},\ B_{\mathrm{emp}}$ & Baseline, theory-preconditioned, and empirically preconditioned sampled operators. \\[0.4em]
$\tau$ & Ridge-regularization parameter for ill-conditioned inverses. \\[0.4em]
$u$ & Unit roundoff (machine precision); the floor $\kappa u$ sets the achievable relative accuracy. \\[0.4em]
$\widehat{r}(t)$ & Sampled residual vector. \\[0.4em]
$\mathbf{1}$ & All-ones vector in $\mathbb{R}^d$. \\[0.6em]

% ----------------------------------------------------------------
\multicolumn{2}{l}{\textbf{Finite-width objects and diagnostics}}\\[0.3em]
$B_m^{(K)}$ & Frozen finite-width tangent block on $\mathcal{H}_K$, $P_K T_0^{(m)} P_K$; $\mathbb{E}[B_m^{(K)}]=\Lambda_K$. \\[0.4em]
$C_t^{(K)}$ & Evolving tangent block on $\mathcal{H}_K$, $P_K T_t^{(m)} P_K$; $C_0^{(K)}=B_m^{(K)}$. \\[0.4em]
$L_t^{(K)}$ & High-to-low frequency coupling (leakage) operator, $P_K T_t^{(m)}(I-P_K)$. \\[0.4em]
$a_t,\ b_t$ & Low- and high-frequency residual components, $a_t=P_K r_t$, $b_t=(I-P_K)r_t$. \\[0.4em]
$\xi_{r,pq}$ & One-neuron matrix entry $\langle\phi_p, T^{(r)}\phi_q\rangle$. \\[0.4em]
$P_K$ & Orthogonal projection $L^2(\mu)\to\mathcal{H}_K$. \\[0.4em]
$P_k,\ U_k$ & Orthogonal projector onto the grid representation of $\mathcal{F}_k$, $P_k=U_kU_k^\top$, with $U_k$ an orthonormal basis. \\[0.4em]
$\widehat{\mathcal{F}}_k(t),\ \widehat P_k(t),\ \widehat U_k(t)$ & Matched empirical eigenspace of the tangent operator, its projector and orthonormal basis. \\[0.4em]
$v_j(t)$ & Orthonormal eigenvectors of the grid-discretized tangent operator. \\[0.4em]
$N,\ \theta_j$ & Number of uniform grid points and grid angles $\theta_j=2\pi j/N$, $j=0,\dots,N-1$. \\[0.4em]
$\varepsilon_{k,F}^{(m)}(t)$ & Projector error $\|\widehat P_k(t)-P_k\|_F$ for frequency $k$ (frozen value at $t=0$ written $\varepsilon_{k,F}^{(m)}$). \\[0.4em]
$\theta_{k,i}(t)$ & $i$-th principal angle between $\widehat{\mathcal{F}}_k(t)$ and $\mathcal{F}_k$; $\cos\theta_{k,i}=s_i(U_k^\top\widehat U_k)$. \\[0.4em]
$s_i(\cdot)$ & $i$-th largest singular value. \\[0.4em]
$s_k^{\mathrm{cos}}(t)$ & Subspace cosine similarity, root-mean-square of the principal-angle cosines for frequency $k$. \\[0.4em]
$\mu_{\mathcal{F}_k}(t)$ & Average spectral strength inside Fourier block $\mathcal{F}_k$, $\tfrac{1}{d_k}\operatorname{tr}(P_k T_t^{(m)} P_k)$. \\[0.4em]
$\mu_{\le K}(t)$ & Cumulative low-frequency spectral strength over the retained block. \\[0.4em]
$\mu_m,\ \mu_t,\ \delta_{\theta_\alpha}$ & Empirical neuron measure, its limit, and the Dirac mass at $\theta_\alpha$ (mean-field discussion). \\[0.6em]

% ----------------------------------------------------------------
\multicolumn{2}{l}{\textbf{Targets and miscellaneous}}\\[0.3em]
$f_{\mathrm{orig}},\ f_{\mathrm{eq}},\ f_{\mathrm{rev}}$ & Original, equal-amplitude, and reversed-amplitude target functions used in the amplitude-variant experiments. \\[0.4em]
$A_i,\ k_i,\ \phi_i$ & Amplitude, frequency, and phase of the $i$-th sinusoid in a synthetic Fourier-regression target. \\[0.4em]
$c$ & Universal constant in the finite-width concentration bound (illustrative value $c=0.2$). \\[0.4em]
$\delta$ & Failure probability in high-probability bounds; also the angular distance $d(\varphi,\psi)$ on $S^1$. \\[0.4em]
$\varepsilon$ & Deviation level in concentration inequalities. \\[0.4em]
$R,\ R_\alpha$ & Rotation matrix on $\mathbb{R}^2$ (rotation by angle $\alpha$). \\

\end{longtable}

%
%  This uses only the longtable environment from the `longtable` package.
%  Add \usepackage{longtable} to the preamble if it is not already loaded
%  (amsmath/amssymb are already loaded in thesis.tex).
% ---------------------------------------------------------------------

\chapter*{Notation}
\addcontentsline{toc}{chapter}{Notation}
\markboth{Notation}{Notation}

The following list summarizes the notation used throughout this thesis. Symbols
are grouped by theme and, within each group, ordered roughly by first
appearance. Page or equation references are given where a symbol is formally
defined. Throughout, the circle angle is written $\theta$ in the introductory
linear-model discussion (Chapter~\ref{cha:intro}) and $\varphi$ from the
methodology onward (Chapters~\ref{cha:methodology}--\ref{cha:results}); the two
denote the same angular coordinate on $S^1$.

\vspace{1em}

\begin{longtable}{p{0.26\textwidth} p{0.68\textwidth}}

% ----------------------------------------------------------------
\multicolumn{2}{l}{\textbf{Sets, spaces, and measures}}\\[0.3em]
\endfirsthead
\multicolumn{2}{l}{\small\itshape (Notation, continued)}\\[0.3em]
\endhead

$\mathbb{R},\ \mathbb{Z}$ & The real numbers and the integers. \\[0.4em]
$\mathbb{R}^{d}$ & $d$-dimensional Euclidean space. \\[0.4em]
$S^1$ & The unit circle, identified with the angular domain $\Omega=[0,2\pi)$ and embedded in $\mathbb{R}^2$ via $x(\varphi)=(\cos\varphi,\sin\varphi)$. \\[0.4em]
$S^{d-1}$ & The unit sphere in $\mathbb{R}^d$ (the $(d-1)$-sphere). \\[0.4em]
$\mathcal{X},\ \mathcal{Y}$ & Input space ($\mathcal{X}\subseteq\mathbb{R}^d$) and target space (e.g.\ $\mathcal{Y}=\mathbb{R}^n$ in regression, $\mathcal{Y}=\{1,\dots,C\}$ in classification). \\[0.4em]
$\rho$ & Joint data distribution over $\mathcal{X}\times\mathcal{Y}$; also used for the uniform input distribution. \\[0.4em]
$\mu$ & Uniform probability measure on $S^1$, $d\mu(\varphi)=d\varphi/(2\pi)$. \\[0.4em]
$L^2(\mu)$ & Space of square-integrable functions on $S^1$ with respect to $\mu$, with inner product $\langle\cdot,\cdot\rangle_{L^2(\mu)}$. \\[0.4em]
$\mathcal{H}$ & A reproducing kernel Hilbert space (RKHS). \\[0.4em]
$\mathcal{H}_\theta$ & Tangent parameter space ($\mathbb{R}^p$ in finite dimension; its limit in the infinite-width case). \\[0.4em]
$\Omega$ & The angular domain $[0,2\pi)$; in the mean-field discussion, the neuron parameter space $\mathbb{R}\times\mathbb{R}^d$. \\[0.4em]
$\|\cdot\|_2,\ \|\cdot\|_F,\ \|\cdot\|_{\max}$ & Euclidean ($\ell^2$) norm, Frobenius norm, and entrywise max norm, respectively. \\[0.4em]
$\langle\cdot,\cdot\rangle$ & Euclidean / Hilbert-space inner product. \\[0.4em]
$\mathbb{E},\ \mathbb{P}$ & Expectation and probability. \\[0.4em]
$\mathcal{N}(0,\sigma^2)$ & Gaussian distribution with mean $0$ and variance $\sigma^2$; $I_d$ denotes the $d\times d$ identity. \\[0.6em]

% ----------------------------------------------------------------
\multicolumn{2}{l}{\textbf{Network, data, and training}}\\[0.3em]
$f(x;\theta),\ f_\theta$ & Neural network output at input $x$ with parameter vector $\theta$. \\[0.4em]
$f_m(x;\theta)$ & One-hidden-layer ReLU network of width $m$ in NTK parameterization, $f_m(x;\theta)=\tfrac{1}{\sqrt{m}}\sum_{i=1}^m a_i\,\sigma(w_i^\top x+b_i)$. \\[0.4em]
$\theta$ & Parameter vector, $\theta\in\mathbb{R}^p$; $\theta_0$ its value at initialization. \\[0.4em]
$p$ & Number of network parameters. \\[0.4em]
$m$ & Network width (number of hidden units); also $w$ as a width label in figures. \\[0.4em]
$d$ & Input dimension; also the dimension $d=2K+1$ of the truncated Fourier block $\mathcal{H}_K$. \\[0.4em]
$a_i,\ w_i,\ b_i$ & Output weight, input weight vector, and bias of hidden unit $i$; initialized $a_i\sim\mathcal{N}(0,1)$, $w_i\sim\mathcal{N}(0,I_2)$, $b_i(0)=0$. \\[0.4em]
$v_\alpha,\ w_\alpha$ & Second-layer and first-layer parameters of neuron $\alpha$ in the generic two-layer network. \\[0.4em]
$\sigma,\ \varphi(\cdot)$ & Activation function; $\sigma(t)=t_+=\max\{t,0\}$ is ReLU ($\varphi$ also denotes a generic activation, and elsewhere the circle angle). \\[0.4em]
$W,\ v$ & Hidden-layer weight matrix and output-weight vector of a two-layer network. \\[0.4em]
$\odot$ & Hadamard (elementwise) product. \\[0.4em]
$\mathcal{D}=\{(x_i,y_i)\}_{i=1}^n$ & Training dataset of $n$ i.i.d.\ samples. \\[0.4em]
$n$ & Number of training samples / sampled points. \\[0.4em]
$X$ & Data matrix ($X\in\mathbb{R}^{d\times n}$, columns $x_i$). \\[0.4em]
$y$ & Vector / function of targets. \\[0.4em]
$\ell(\cdot,\cdot)$ & Loss function. \\[0.4em]
$\mathcal{L},\ L$ & Empirical (least-squares) risk; $\mathcal{L}_\rho$ denotes the population risk. \\[0.4em]
$\eta$ & Gradient-descent step size / learning rate. \\[0.4em]
$t,\ k$ (subscript) & Training time / iteration index. \\[0.4em]
$\nabla_\theta$ & Gradient with respect to $\theta$; $\dot{(\cdot)}$ denotes time derivative under gradient flow. \\[0.4em]
$\theta^\star,\ w^\star,\ f^\star$ & Minimizers / fixed points of the relevant objective. \\[0.4em]
$A^{+},\ A^{\ast}$ & Moore--Penrose pseudoinverse and adjoint of an operator/matrix $A$. \\[0.4em]
$\mathrm{range}(\cdot),\ \ker(\cdot)$ & Range (column space) and kernel (null space) of an operator. \\[0.4em]
$\lambda_{\max}(\cdot)$ & Largest eigenvalue of a matrix. \\[0.6em]

% ----------------------------------------------------------------
\multicolumn{2}{l}{\textbf{Neural Tangent Kernel and operators}}\\[0.3em]
$\phi(x),\ \Phi_t(x)$ & Tangent feature map $\nabla_\theta f(x;\theta_0)$ (and its time-$t$ version). \\[0.4em]
$\Theta_0(x,x')$ & Neural Tangent Kernel at initialization, $\nabla_\theta f(x;\theta_0)^\top\nabla_\theta f(x';\theta_0)$. \\[0.4em]
$\Theta_t,\ \Theta_t^{(m)}$ & NTK at training time $t$ (the empirical, width-$m$ kernel). \\[0.4em]
$\Theta_0^{(m)}$ & Finite-width empirical NTK at initialization. \\[0.4em]
$\bar\Theta,\ \Theta$ & Limiting (infinite-width) deterministic NTK; also written as the Gram matrix $\bar\Theta=(\bar\Theta(x_i,x_j))_{i,j}$. \\[0.4em]
$\Theta(\delta)$ & Limiting NTK on $S^1$ as a function of the angular distance $\delta=d(\varphi,\psi)\in[0,\pi]$. \\[0.4em]
$\Theta^{(\mathrm{nb})}$ & Limiting NTK with the trainable hidden-bias contribution removed. \\[0.4em]
$\Theta^{(r)},\ \psi_r$ & One-neuron kernel contribution and one-neuron tangent feature of neuron $r$. \\[0.4em]
$J_0$ & Jacobian of network outputs w.r.t.\ parameters at initialization; $\bar\Theta=J_0 J_0^\ast$. \\[0.4em]
$T_{\bar\Theta},\ T,\ T_\infty$ & Integral operator on $L^2(\mu)$ induced by the limiting kernel; $T_\infty=\lim_{m\to\infty}T_0^{(m)}$. \\[0.4em]
$T_t^{(m)},\ T_0^{(m)}$ & Empirical tangent integral operator at time $t$ and at initialization (frozen operator). \\[0.4em]
$T^{(r)}$ & One-neuron integral operator; $T_0^{(m)}=\tfrac1m\sum_{r=1}^m T^{(r)}$. \\[0.4em]
$T_p$ & Density-weighted integral operator for a non-uniform input density $p(\theta)$. \\[0.4em]
$\kappa$ & The $2\pi$-periodic kernel profile with $\bar\Theta(\theta,\theta')=\kappa(\theta-\theta')$; also a uniform kernel bound and (in the conditioning discussion) a matrix condition number. \\[0.4em]
$r_t,\ r$ & Residual $f_t-y$ (function or vector). \\[0.6em]

% ----------------------------------------------------------------
\multicolumn{2}{l}{\textbf{Fourier basis, spectrum, and subspaces}}\\[0.3em]
$\theta,\ \varphi,\ \gamma$ & Angular coordinate on $S^1$; $\gamma$ is used as the angle variable in target functions. \\[0.4em]
$\phi_q(\theta)=e^{iq\theta}$ & Complex Fourier mode of frequency $q\in\mathbb{Z}$. \\[0.4em]
$\phi_0,\ \phi_{k,c},\ \phi_{k,s}$ & Real Fourier basis of $L^2(\mu)$: $\phi_0=1$, $\phi_{k,c}=\sqrt{2}\cos(k\varphi)$, $\phi_{k,s}=\sqrt{2}\sin(k\varphi)$. \\[0.4em]
$\psi_j$ & Generic orthonormal eigenfunction of an integral operator. \\[0.4em]
$k,\ q$ & Fourier frequency index; $k(p)$ is the frequency of basis function $\phi_p$. \\[0.4em]
$K$ & Frequency cutoff defining the retained low-frequency block. \\[0.4em]
$\mathcal{H}_K$ & Truncated Fourier subspace $\operatorname{span}\{\phi_0,\phi_{k,c},\phi_{k,s}:1\le k\le K\}$, of dimension $d=2K+1$. \\[0.4em]
$\mathcal{F}_k$ & Frequency-$k$ subspace: $\mathcal{F}_0=\operatorname{span}\{\phi_0\}$, $\mathcal{F}_k=\operatorname{span}\{\phi_{k,c},\phi_{k,s}\}$ for $k\ge1$. \\[0.4em]
$d_k$ & Dimension of $\mathcal{F}_k$: $d_0=1$, $d_k=2$ for $k\ge1$. \\[0.4em]
$\lambda_q,\ \lambda_k$ & Fourier eigenvalue of $T$ at frequency $q$ (resp.\ $k$); $\lambda_p:=\lambda_{k(p)}$. \\[0.4em]
$\lambda_k^{(\mathrm{nb})}$ & Fourier eigenvalue of the no-bias kernel $\Theta^{(\mathrm{nb})}$. \\[0.4em]
$\Lambda$ & Diagonal matrix of eigenvalues $\operatorname{diag}(\lambda_p)$; in the linear model, $A=Q\Lambda Q^\top$. \\[0.4em]
$\Lambda_K$ & Continuum reference block $P_K T_\infty P_K=\operatorname{diag}(\lambda_0,\lambda_1,\lambda_1,\dots,\lambda_K,\lambda_K)$. \\[0.4em]
$\alpha_p(t)$ & Amplitude of the residual along Fourier mode $\phi_p$; decays as $e^{-\lambda_p t}\alpha_p(0)$. \\[0.4em]
$a_j(t)$ & Coefficient of the residual along eigenfunction $\psi_j$ in the continuum operator analysis. \\[0.6em]

% ----------------------------------------------------------------
\multicolumn{2}{l}{\textbf{Finite-sample objects}}\\[0.3em]
$\varphi_1,\dots,\varphi_n$ & I.i.d.\ angles sampled uniformly on $[0,2\pi)$. \\[0.4em]
$\Phi$ & Sampled Fourier feature matrix, $\Phi\in\mathbb{R}^{n\times d}$, $\Phi_{i,p}=\phi_p(\varphi_i)$; $\Phi_p$ its $p$-th column. \\[0.4em]
$A$ & Normalized sampled kernel matrix, $A_{ij}=\tfrac1n\Theta(\varphi_i-\varphi_j)$. \\[0.4em]
$G$ & Sampled Fourier Gram matrix, $G=\tfrac1n\Phi^\top\Phi$. \\[0.4em]
$H$ & Compressed operator matrix, $H=\tfrac1n\Phi^\top A\Phi$; $\mathbb{E}[H]=\Lambda^{(n)}$. \\[0.4em]
$\lambda_p^{(n)}$ & Finite-sample corrected Fourier eigenvalue, $(1-\tfrac1n)\lambda_p+\tfrac{\Theta(0)}{n}$. \\[0.4em]
$\Lambda^{(n)}$ & Diagonal matrix $\operatorname{diag}(\lambda_p^{(n)})$ of corrected eigenvalues. \\[0.4em]
$C,\ C_n^{(K)}$ & Coefficient-level restricted operator $G^{-1}H$ on the retained block; compared against $\Lambda^{(n)}$. \\[0.4em]
$P_{\mathrm{th}},\ P_{\mathrm{emp}}$ & Theory and empirical Fourier-block preconditioners. \\[0.4em]
$B_{\mathrm{base}},\ B_{\mathrm{th}},\ B_{\mathrm{emp}}$ & Baseline, theory-preconditioned, and empirically preconditioned sampled operators. \\[0.4em]
$\tau$ & Ridge-regularization parameter for ill-conditioned inverses. \\[0.4em]
$u$ & Unit roundoff (machine precision); the floor $\kappa u$ sets the achievable relative accuracy. \\[0.4em]
$\widehat{r}(t)$ & Sampled residual vector. \\[0.4em]
$\mathbf{1}$ & All-ones vector in $\mathbb{R}^d$. \\[0.6em]

% ----------------------------------------------------------------
\multicolumn{2}{l}{\textbf{Finite-width objects and diagnostics}}\\[0.3em]
$B_m^{(K)}$ & Frozen finite-width tangent block on $\mathcal{H}_K$, $P_K T_0^{(m)} P_K$; $\mathbb{E}[B_m^{(K)}]=\Lambda_K$. \\[0.4em]
$C_t^{(K)}$ & Evolving tangent block on $\mathcal{H}_K$, $P_K T_t^{(m)} P_K$; $C_0^{(K)}=B_m^{(K)}$. \\[0.4em]
$L_t^{(K)}$ & High-to-low frequency coupling (leakage) operator, $P_K T_t^{(m)}(I-P_K)$. \\[0.4em]
$a_t,\ b_t$ & Low- and high-frequency residual components, $a_t=P_K r_t$, $b_t=(I-P_K)r_t$. \\[0.4em]
$\xi_{r,pq}$ & One-neuron matrix entry $\langle\phi_p, T^{(r)}\phi_q\rangle$. \\[0.4em]
$P_K$ & Orthogonal projection $L^2(\mu)\to\mathcal{H}_K$. \\[0.4em]
$P_k,\ U_k$ & Orthogonal projector onto the grid representation of $\mathcal{F}_k$, $P_k=U_kU_k^\top$, with $U_k$ an orthonormal basis. \\[0.4em]
$\widehat{\mathcal{F}}_k(t),\ \widehat P_k(t),\ \widehat U_k(t)$ & Matched empirical eigenspace of the tangent operator, its projector and orthonormal basis. \\[0.4em]
$v_j(t)$ & Orthonormal eigenvectors of the grid-discretized tangent operator. \\[0.4em]
$N,\ \theta_j$ & Number of uniform grid points and grid angles $\theta_j=2\pi j/N$, $j=0,\dots,N-1$. \\[0.4em]
$\varepsilon_{k,F}^{(m)}(t)$ & Projector error $\|\widehat P_k(t)-P_k\|_F$ for frequency $k$ (frozen value at $t=0$ written $\varepsilon_{k,F}^{(m)}$). \\[0.4em]
$\theta_{k,i}(t)$ & $i$-th principal angle between $\widehat{\mathcal{F}}_k(t)$ and $\mathcal{F}_k$; $\cos\theta_{k,i}=s_i(U_k^\top\widehat U_k)$. \\[0.4em]
$s_i(\cdot)$ & $i$-th largest singular value. \\[0.4em]
$s_k^{\mathrm{cos}}(t)$ & Subspace cosine similarity, root-mean-square of the principal-angle cosines for frequency $k$. \\[0.4em]
$\mu_{\mathcal{F}_k}(t)$ & Average spectral strength inside Fourier block $\mathcal{F}_k$, $\tfrac{1}{d_k}\operatorname{tr}(P_k T_t^{(m)} P_k)$. \\[0.4em]
$\mu_{\le K}(t)$ & Cumulative low-frequency spectral strength over the retained block. \\[0.4em]
$\mu_m,\ \mu_t,\ \delta_{\theta_\alpha}$ & Empirical neuron measure, its limit, and the Dirac mass at $\theta_\alpha$ (mean-field discussion). \\[0.6em]

% ----------------------------------------------------------------
\multicolumn{2}{l}{\textbf{Targets and miscellaneous}}\\[0.3em]
$f_{\mathrm{orig}},\ f_{\mathrm{eq}},\ f_{\mathrm{rev}}$ & Original, equal-amplitude, and reversed-amplitude target functions used in the amplitude-variant experiments. \\[0.4em]
$A_i,\ k_i,\ \phi_i$ & Amplitude, frequency, and phase of the $i$-th sinusoid in a synthetic Fourier-regression target. \\[0.4em]
$c$ & Universal constant in the finite-width concentration bound (illustrative value $c=0.2$). \\[0.4em]
$\delta$ & Failure probability in high-probability bounds; also the angular distance $d(\varphi,\psi)$ on $S^1$. \\[0.4em]
$\varepsilon$ & Deviation level in concentration inequalities. \\[0.4em]
$R,\ R_\alpha$ & Rotation matrix on $\mathbb{R}^2$ (rotation by angle $\alpha$). \\

\end{longtable}

%
%  This uses only the longtable environment from the `longtable` package.
%  Add \usepackage{longtable} to the preamble if it is not already loaded
%  (amsmath/amssymb are already loaded in thesis.tex).
% ---------------------------------------------------------------------

\chapter*{Notation}
\addcontentsline{toc}{chapter}{Notation}
\markboth{Notation}{Notation}

The following list summarizes the notation used throughout this thesis. Symbols
are grouped by theme and, within each group, ordered roughly by first
appearance. Page or equation references are given where a symbol is formally
defined. Throughout, the circle angle is written $\theta$ in the introductory
linear-model discussion (Chapter~\ref{cha:intro}) and $\varphi$ from the
methodology onward (Chapters~\ref{cha:methodology}--\ref{cha:results}); the two
denote the same angular coordinate on $S^1$.

\vspace{1em}

\begin{longtable}{p{0.26\textwidth} p{0.68\textwidth}}

% ----------------------------------------------------------------
\multicolumn{2}{l}{\textbf{Sets, spaces, and measures}}\\[0.3em]
\endfirsthead
\multicolumn{2}{l}{\small\itshape (Notation, continued)}\\[0.3em]
\endhead

$\mathbb{R},\ \mathbb{Z}$ & The real numbers and the integers. \\[0.4em]
$\mathbb{R}^{d}$ & $d$-dimensional Euclidean space. \\[0.4em]
$S^1$ & The unit circle, identified with the angular domain $\Omega=[0,2\pi)$ and embedded in $\mathbb{R}^2$ via $x(\varphi)=(\cos\varphi,\sin\varphi)$. \\[0.4em]
$S^{d-1}$ & The unit sphere in $\mathbb{R}^d$ (the $(d-1)$-sphere). \\[0.4em]
$\mathcal{X},\ \mathcal{Y}$ & Input space ($\mathcal{X}\subseteq\mathbb{R}^d$) and target space (e.g.\ $\mathcal{Y}=\mathbb{R}^n$ in regression, $\mathcal{Y}=\{1,\dots,C\}$ in classification). \\[0.4em]
$\rho$ & Joint data distribution over $\mathcal{X}\times\mathcal{Y}$; also used for the uniform input distribution. \\[0.4em]
$\mu$ & Uniform probability measure on $S^1$, $d\mu(\varphi)=d\varphi/(2\pi)$. \\[0.4em]
$L^2(\mu)$ & Space of square-integrable functions on $S^1$ with respect to $\mu$, with inner product $\langle\cdot,\cdot\rangle_{L^2(\mu)}$. \\[0.4em]
$\mathcal{H}$ & A reproducing kernel Hilbert space (RKHS). \\[0.4em]
$\mathcal{H}_\theta$ & Tangent parameter space ($\mathbb{R}^p$ in finite dimension; its limit in the infinite-width case). \\[0.4em]
$\Omega$ & The angular domain $[0,2\pi)$; in the mean-field discussion, the neuron parameter space $\mathbb{R}\times\mathbb{R}^d$. \\[0.4em]
$\|\cdot\|_2,\ \|\cdot\|_F,\ \|\cdot\|_{\max}$ & Euclidean ($\ell^2$) norm, Frobenius norm, and entrywise max norm, respectively. \\[0.4em]
$\langle\cdot,\cdot\rangle$ & Euclidean / Hilbert-space inner product. \\[0.4em]
$\mathbb{E},\ \mathbb{P}$ & Expectation and probability. \\[0.4em]
$\mathcal{N}(0,\sigma^2)$ & Gaussian distribution with mean $0$ and variance $\sigma^2$; $I_d$ denotes the $d\times d$ identity. \\[0.6em]

% ----------------------------------------------------------------
\multicolumn{2}{l}{\textbf{Network, data, and training}}\\[0.3em]
$f(x;\theta),\ f_\theta$ & Neural network output at input $x$ with parameter vector $\theta$. \\[0.4em]
$f_m(x;\theta)$ & One-hidden-layer ReLU network of width $m$ in NTK parameterization, $f_m(x;\theta)=\tfrac{1}{\sqrt{m}}\sum_{i=1}^m a_i\,\sigma(w_i^\top x+b_i)$. \\[0.4em]
$\theta$ & Parameter vector, $\theta\in\mathbb{R}^p$; $\theta_0$ its value at initialization. \\[0.4em]
$p$ & Number of network parameters. \\[0.4em]
$m$ & Network width (number of hidden units); also $w$ as a width label in figures. \\[0.4em]
$d$ & Input dimension; also the dimension $d=2K+1$ of the truncated Fourier block $\mathcal{H}_K$. \\[0.4em]
$a_i,\ w_i,\ b_i$ & Output weight, input weight vector, and bias of hidden unit $i$; initialized $a_i\sim\mathcal{N}(0,1)$, $w_i\sim\mathcal{N}(0,I_2)$, $b_i(0)=0$. \\[0.4em]
$v_\alpha,\ w_\alpha$ & Second-layer and first-layer parameters of neuron $\alpha$ in the generic two-layer network. \\[0.4em]
$\sigma,\ \varphi(\cdot)$ & Activation function; $\sigma(t)=t_+=\max\{t,0\}$ is ReLU ($\varphi$ also denotes a generic activation, and elsewhere the circle angle). \\[0.4em]
$W,\ v$ & Hidden-layer weight matrix and output-weight vector of a two-layer network. \\[0.4em]
$\odot$ & Hadamard (elementwise) product. \\[0.4em]
$\mathcal{D}=\{(x_i,y_i)\}_{i=1}^n$ & Training dataset of $n$ i.i.d.\ samples. \\[0.4em]
$n$ & Number of training samples / sampled points. \\[0.4em]
$X$ & Data matrix ($X\in\mathbb{R}^{d\times n}$, columns $x_i$). \\[0.4em]
$y$ & Vector / function of targets. \\[0.4em]
$\ell(\cdot,\cdot)$ & Loss function. \\[0.4em]
$\mathcal{L},\ L$ & Empirical (least-squares) risk; $\mathcal{L}_\rho$ denotes the population risk. \\[0.4em]
$\eta$ & Gradient-descent step size / learning rate. \\[0.4em]
$t,\ k$ (subscript) & Training time / iteration index. \\[0.4em]
$\nabla_\theta$ & Gradient with respect to $\theta$; $\dot{(\cdot)}$ denotes time derivative under gradient flow. \\[0.4em]
$\theta^\star,\ w^\star,\ f^\star$ & Minimizers / fixed points of the relevant objective. \\[0.4em]
$A^{+},\ A^{\ast}$ & Moore--Penrose pseudoinverse and adjoint of an operator/matrix $A$. \\[0.4em]
$\mathrm{range}(\cdot),\ \ker(\cdot)$ & Range (column space) and kernel (null space) of an operator. \\[0.4em]
$\lambda_{\max}(\cdot)$ & Largest eigenvalue of a matrix. \\[0.6em]

% ----------------------------------------------------------------
\multicolumn{2}{l}{\textbf{Neural Tangent Kernel and operators}}\\[0.3em]
$\phi(x),\ \Phi_t(x)$ & Tangent feature map $\nabla_\theta f(x;\theta_0)$ (and its time-$t$ version). \\[0.4em]
$\Theta_0(x,x')$ & Neural Tangent Kernel at initialization, $\nabla_\theta f(x;\theta_0)^\top\nabla_\theta f(x';\theta_0)$. \\[0.4em]
$\Theta_t,\ \Theta_t^{(m)}$ & NTK at training time $t$ (the empirical, width-$m$ kernel). \\[0.4em]
$\Theta_0^{(m)}$ & Finite-width empirical NTK at initialization. \\[0.4em]
$\bar\Theta,\ \Theta$ & Limiting (infinite-width) deterministic NTK; also written as the Gram matrix $\bar\Theta=(\bar\Theta(x_i,x_j))_{i,j}$. \\[0.4em]
$\Theta(\delta)$ & Limiting NTK on $S^1$ as a function of the angular distance $\delta=d(\varphi,\psi)\in[0,\pi]$. \\[0.4em]
$\Theta^{(\mathrm{nb})}$ & Limiting NTK with the trainable hidden-bias contribution removed. \\[0.4em]
$\Theta^{(r)},\ \psi_r$ & One-neuron kernel contribution and one-neuron tangent feature of neuron $r$. \\[0.4em]
$J_0$ & Jacobian of network outputs w.r.t.\ parameters at initialization; $\bar\Theta=J_0 J_0^\ast$. \\[0.4em]
$T_{\bar\Theta},\ T,\ T_\infty$ & Integral operator on $L^2(\mu)$ induced by the limiting kernel; $T_\infty=\lim_{m\to\infty}T_0^{(m)}$. \\[0.4em]
$T_t^{(m)},\ T_0^{(m)}$ & Empirical tangent integral operator at time $t$ and at initialization (frozen operator). \\[0.4em]
$T^{(r)}$ & One-neuron integral operator; $T_0^{(m)}=\tfrac1m\sum_{r=1}^m T^{(r)}$. \\[0.4em]
$T_p$ & Density-weighted integral operator for a non-uniform input density $p(\theta)$. \\[0.4em]
$\kappa$ & The $2\pi$-periodic kernel profile with $\bar\Theta(\theta,\theta')=\kappa(\theta-\theta')$; also a uniform kernel bound and (in the conditioning discussion) a matrix condition number. \\[0.4em]
$r_t,\ r$ & Residual $f_t-y$ (function or vector). \\[0.6em]

% ----------------------------------------------------------------
\multicolumn{2}{l}{\textbf{Fourier basis, spectrum, and subspaces}}\\[0.3em]
$\theta,\ \varphi,\ \gamma$ & Angular coordinate on $S^1$; $\gamma$ is used as the angle variable in target functions. \\[0.4em]
$\phi_q(\theta)=e^{iq\theta}$ & Complex Fourier mode of frequency $q\in\mathbb{Z}$. \\[0.4em]
$\phi_0,\ \phi_{k,c},\ \phi_{k,s}$ & Real Fourier basis of $L^2(\mu)$: $\phi_0=1$, $\phi_{k,c}=\sqrt{2}\cos(k\varphi)$, $\phi_{k,s}=\sqrt{2}\sin(k\varphi)$. \\[0.4em]
$\psi_j$ & Generic orthonormal eigenfunction of an integral operator. \\[0.4em]
$k,\ q$ & Fourier frequency index; $k(p)$ is the frequency of basis function $\phi_p$. \\[0.4em]
$K$ & Frequency cutoff defining the retained low-frequency block. \\[0.4em]
$\mathcal{H}_K$ & Truncated Fourier subspace $\operatorname{span}\{\phi_0,\phi_{k,c},\phi_{k,s}:1\le k\le K\}$, of dimension $d=2K+1$. \\[0.4em]
$\mathcal{F}_k$ & Frequency-$k$ subspace: $\mathcal{F}_0=\operatorname{span}\{\phi_0\}$, $\mathcal{F}_k=\operatorname{span}\{\phi_{k,c},\phi_{k,s}\}$ for $k\ge1$. \\[0.4em]
$d_k$ & Dimension of $\mathcal{F}_k$: $d_0=1$, $d_k=2$ for $k\ge1$. \\[0.4em]
$\lambda_q,\ \lambda_k$ & Fourier eigenvalue of $T$ at frequency $q$ (resp.\ $k$); $\lambda_p:=\lambda_{k(p)}$. \\[0.4em]
$\lambda_k^{(\mathrm{nb})}$ & Fourier eigenvalue of the no-bias kernel $\Theta^{(\mathrm{nb})}$. \\[0.4em]
$\Lambda$ & Diagonal matrix of eigenvalues $\operatorname{diag}(\lambda_p)$; in the linear model, $A=Q\Lambda Q^\top$. \\[0.4em]
$\Lambda_K$ & Continuum reference block $P_K T_\infty P_K=\operatorname{diag}(\lambda_0,\lambda_1,\lambda_1,\dots,\lambda_K,\lambda_K)$. \\[0.4em]
$\alpha_p(t)$ & Amplitude of the residual along Fourier mode $\phi_p$; decays as $e^{-\lambda_p t}\alpha_p(0)$. \\[0.4em]
$a_j(t)$ & Coefficient of the residual along eigenfunction $\psi_j$ in the continuum operator analysis. \\[0.6em]

% ----------------------------------------------------------------
\multicolumn{2}{l}{\textbf{Finite-sample objects}}\\[0.3em]
$\varphi_1,\dots,\varphi_n$ & I.i.d.\ angles sampled uniformly on $[0,2\pi)$. \\[0.4em]
$\Phi$ & Sampled Fourier feature matrix, $\Phi\in\mathbb{R}^{n\times d}$, $\Phi_{i,p}=\phi_p(\varphi_i)$; $\Phi_p$ its $p$-th column. \\[0.4em]
$A$ & Normalized sampled kernel matrix, $A_{ij}=\tfrac1n\Theta(\varphi_i-\varphi_j)$. \\[0.4em]
$G$ & Sampled Fourier Gram matrix, $G=\tfrac1n\Phi^\top\Phi$. \\[0.4em]
$H$ & Compressed operator matrix, $H=\tfrac1n\Phi^\top A\Phi$; $\mathbb{E}[H]=\Lambda^{(n)}$. \\[0.4em]
$\lambda_p^{(n)}$ & Finite-sample corrected Fourier eigenvalue, $(1-\tfrac1n)\lambda_p+\tfrac{\Theta(0)}{n}$. \\[0.4em]
$\Lambda^{(n)}$ & Diagonal matrix $\operatorname{diag}(\lambda_p^{(n)})$ of corrected eigenvalues. \\[0.4em]
$C,\ C_n^{(K)}$ & Coefficient-level restricted operator $G^{-1}H$ on the retained block; compared against $\Lambda^{(n)}$. \\[0.4em]
$P_{\mathrm{th}},\ P_{\mathrm{emp}}$ & Theory and empirical Fourier-block preconditioners. \\[0.4em]
$B_{\mathrm{base}},\ B_{\mathrm{th}},\ B_{\mathrm{emp}}$ & Baseline, theory-preconditioned, and empirically preconditioned sampled operators. \\[0.4em]
$\tau$ & Ridge-regularization parameter for ill-conditioned inverses. \\[0.4em]
$u$ & Unit roundoff (machine precision); the floor $\kappa u$ sets the achievable relative accuracy. \\[0.4em]
$\widehat{r}(t)$ & Sampled residual vector. \\[0.4em]
$\mathbf{1}$ & All-ones vector in $\mathbb{R}^d$. \\[0.6em]

% ----------------------------------------------------------------
\multicolumn{2}{l}{\textbf{Finite-width objects and diagnostics}}\\[0.3em]
$B_m^{(K)}$ & Frozen finite-width tangent block on $\mathcal{H}_K$, $P_K T_0^{(m)} P_K$; $\mathbb{E}[B_m^{(K)}]=\Lambda_K$. \\[0.4em]
$C_t^{(K)}$ & Evolving tangent block on $\mathcal{H}_K$, $P_K T_t^{(m)} P_K$; $C_0^{(K)}=B_m^{(K)}$. \\[0.4em]
$L_t^{(K)}$ & High-to-low frequency coupling (leakage) operator, $P_K T_t^{(m)}(I-P_K)$. \\[0.4em]
$a_t,\ b_t$ & Low- and high-frequency residual components, $a_t=P_K r_t$, $b_t=(I-P_K)r_t$. \\[0.4em]
$\xi_{r,pq}$ & One-neuron matrix entry $\langle\phi_p, T^{(r)}\phi_q\rangle$. \\[0.4em]
$P_K$ & Orthogonal projection $L^2(\mu)\to\mathcal{H}_K$. \\[0.4em]
$P_k,\ U_k$ & Orthogonal projector onto the grid representation of $\mathcal{F}_k$, $P_k=U_kU_k^\top$, with $U_k$ an orthonormal basis. \\[0.4em]
$\widehat{\mathcal{F}}_k(t),\ \widehat P_k(t),\ \widehat U_k(t)$ & Matched empirical eigenspace of the tangent operator, its projector and orthonormal basis. \\[0.4em]
$v_j(t)$ & Orthonormal eigenvectors of the grid-discretized tangent operator. \\[0.4em]
$N,\ \theta_j$ & Number of uniform grid points and grid angles $\theta_j=2\pi j/N$, $j=0,\dots,N-1$. \\[0.4em]
$\varepsilon_{k,F}^{(m)}(t)$ & Projector error $\|\widehat P_k(t)-P_k\|_F$ for frequency $k$ (frozen value at $t=0$ written $\varepsilon_{k,F}^{(m)}$). \\[0.4em]
$\theta_{k,i}(t)$ & $i$-th principal angle between $\widehat{\mathcal{F}}_k(t)$ and $\mathcal{F}_k$; $\cos\theta_{k,i}=s_i(U_k^\top\widehat U_k)$. \\[0.4em]
$s_i(\cdot)$ & $i$-th largest singular value. \\[0.4em]
$s_k^{\mathrm{cos}}(t)$ & Subspace cosine similarity, root-mean-square of the principal-angle cosines for frequency $k$. \\[0.4em]
$\mu_{\mathcal{F}_k}(t)$ & Average spectral strength inside Fourier block $\mathcal{F}_k$, $\tfrac{1}{d_k}\operatorname{tr}(P_k T_t^{(m)} P_k)$. \\[0.4em]
$\mu_{\le K}(t)$ & Cumulative low-frequency spectral strength over the retained block. \\[0.4em]
$\mu_m,\ \mu_t,\ \delta_{\theta_\alpha}$ & Empirical neuron measure, its limit, and the Dirac mass at $\theta_\alpha$ (mean-field discussion). \\[0.6em]

% ----------------------------------------------------------------
\multicolumn{2}{l}{\textbf{Targets and miscellaneous}}\\[0.3em]
$f_{\mathrm{orig}},\ f_{\mathrm{eq}},\ f_{\mathrm{rev}}$ & Original, equal-amplitude, and reversed-amplitude target functions used in the amplitude-variant experiments. \\[0.4em]
$A_i,\ k_i,\ \phi_i$ & Amplitude, frequency, and phase of the $i$-th sinusoid in a synthetic Fourier-regression target. \\[0.4em]
$c$ & Universal constant in the finite-width concentration bound (illustrative value $c=0.2$). \\[0.4em]
$\delta$ & Failure probability in high-probability bounds; also the angular distance $d(\varphi,\psi)$ on $S^1$. \\[0.4em]
$\varepsilon$ & Deviation level in concentration inequalities. \\[0.4em]
$R,\ R_\alpha$ & Rotation matrix on $\mathbb{R}^2$ (rotation by angle $\alpha$). \\

\end{longtable}

%
%  This uses only the longtable environment from the `longtable` package.
%  Add \usepackage{longtable} to the preamble if it is not already loaded
%  (amsmath/amssymb are already loaded in thesis.tex).
% ---------------------------------------------------------------------

\chapter*{Notation}
\addcontentsline{toc}{chapter}{Notation}
\markboth{Notation}{Notation}

The following list summarizes the notation used throughout this thesis. Symbols
are grouped by theme and, within each group, ordered roughly by first
appearance. Page or equation references are given where a symbol is formally
defined. Throughout, the circle angle is written $\theta$ in the introductory
linear-model discussion (Chapter~\ref{cha:intro}) and $\varphi$ from the
methodology onward (Chapters~\ref{cha:methodology}--\ref{cha:results}); the two
denote the same angular coordinate on $S^1$.

\vspace{1em}

\begin{longtable}{p{0.26\textwidth} p{0.68\textwidth}}

% ----------------------------------------------------------------
\multicolumn{2}{l}{\textbf{Sets, spaces, and measures}}\\[0.3em]
\endfirsthead
\multicolumn{2}{l}{\small\itshape (Notation, continued)}\\[0.3em]
\endhead

$\mathbb{R},\ \mathbb{Z}$ & The real numbers and the integers. \\[0.4em]
$\mathbb{R}^{d}$ & $d$-dimensional Euclidean space. \\[0.4em]
$S^1$ & The unit circle, identified with the angular domain $\Omega=[0,2\pi)$ and embedded in $\mathbb{R}^2$ via $x(\varphi)=(\cos\varphi,\sin\varphi)$. \\[0.4em]
$S^{d-1}$ & The unit sphere in $\mathbb{R}^d$ (the $(d-1)$-sphere). \\[0.4em]
$\mathcal{X},\ \mathcal{Y}$ & Input space ($\mathcal{X}\subseteq\mathbb{R}^d$) and target space (e.g.\ $\mathcal{Y}=\mathbb{R}^n$ in regression, $\mathcal{Y}=\{1,\dots,C\}$ in classification). \\[0.4em]
$\rho$ & Joint data distribution over $\mathcal{X}\times\mathcal{Y}$; also used for the uniform input distribution. \\[0.4em]
$\mu$ & Uniform probability measure on $S^1$, $d\mu(\varphi)=d\varphi/(2\pi)$. \\[0.4em]
$L^2(\mu)$ & Space of square-integrable functions on $S^1$ with respect to $\mu$, with inner product $\langle\cdot,\cdot\rangle_{L^2(\mu)}$. \\[0.4em]
$\mathcal{H}$ & A reproducing kernel Hilbert space (RKHS). \\[0.4em]
$\mathcal{H}_\theta$ & Tangent parameter space ($\mathbb{R}^p$ in finite dimension; its limit in the infinite-width case). \\[0.4em]
$\Omega$ & The angular domain $[0,2\pi)$; in the mean-field discussion, the neuron parameter space $\mathbb{R}\times\mathbb{R}^d$. \\[0.4em]
$\|\cdot\|_2,\ \|\cdot\|_F,\ \|\cdot\|_{\max}$ & Euclidean ($\ell^2$) norm, Frobenius norm, and entrywise max norm, respectively. \\[0.4em]
$\langle\cdot,\cdot\rangle$ & Euclidean / Hilbert-space inner product. \\[0.4em]
$\mathbb{E},\ \mathbb{P}$ & Expectation and probability. \\[0.4em]
$\mathcal{N}(0,\sigma^2)$ & Gaussian distribution with mean $0$ and variance $\sigma^2$; $I_d$ denotes the $d\times d$ identity. \\[0.6em]

% ----------------------------------------------------------------
\multicolumn{2}{l}{\textbf{Network, data, and training}}\\[0.3em]
$f(x;\theta),\ f_\theta$ & Neural network output at input $x$ with parameter vector $\theta$. \\[0.4em]
$f_m(x;\theta)$ & One-hidden-layer ReLU network of width $m$ in NTK parameterization, $f_m(x;\theta)=\tfrac{1}{\sqrt{m}}\sum_{i=1}^m a_i\,\sigma(w_i^\top x+b_i)$. \\[0.4em]
$\theta$ & Parameter vector, $\theta\in\mathbb{R}^p$; $\theta_0$ its value at initialization. \\[0.4em]
$p$ & Number of network parameters. \\[0.4em]
$m$ & Network width (number of hidden units); also $w$ as a width label in figures. \\[0.4em]
$d$ & Input dimension; also the dimension $d=2K+1$ of the truncated Fourier block $\mathcal{H}_K$. \\[0.4em]
$a_i,\ w_i,\ b_i$ & Output weight, input weight vector, and bias of hidden unit $i$; initialized $a_i\sim\mathcal{N}(0,1)$, $w_i\sim\mathcal{N}(0,I_2)$, $b_i(0)=0$. \\[0.4em]
$v_\alpha,\ w_\alpha$ & Second-layer and first-layer parameters of neuron $\alpha$ in the generic two-layer network. \\[0.4em]
$\sigma,\ \varphi(\cdot)$ & Activation function; $\sigma(t)=t_+=\max\{t,0\}$ is ReLU ($\varphi$ also denotes a generic activation, and elsewhere the circle angle). \\[0.4em]
$W,\ v$ & Hidden-layer weight matrix and output-weight vector of a two-layer network. \\[0.4em]
$\odot$ & Hadamard (elementwise) product. \\[0.4em]
$\mathcal{D}=\{(x_i,y_i)\}_{i=1}^n$ & Training dataset of $n$ i.i.d.\ samples. \\[0.4em]
$n$ & Number of training samples / sampled points. \\[0.4em]
$X$ & Data matrix ($X\in\mathbb{R}^{d\times n}$, columns $x_i$). \\[0.4em]
$y$ & Vector / function of targets. \\[0.4em]
$\ell(\cdot,\cdot)$ & Loss function. \\[0.4em]
$\mathcal{L},\ L$ & Empirical (least-squares) risk; $\mathcal{L}_\rho$ denotes the population risk. \\[0.4em]
$\eta$ & Gradient-descent step size / learning rate. \\[0.4em]
$t,\ k$ (subscript) & Training time / iteration index. \\[0.4em]
$\nabla_\theta$ & Gradient with respect to $\theta$; $\dot{(\cdot)}$ denotes time derivative under gradient flow. \\[0.4em]
$\theta^\star,\ w^\star,\ f^\star$ & Minimizers / fixed points of the relevant objective. \\[0.4em]
$A^{+},\ A^{\ast}$ & Moore--Penrose pseudoinverse and adjoint of an operator/matrix $A$. \\[0.4em]
$\mathrm{range}(\cdot),\ \ker(\cdot)$ & Range (column space) and kernel (null space) of an operator. \\[0.4em]
$\lambda_{\max}(\cdot)$ & Largest eigenvalue of a matrix. \\[0.6em]

% ----------------------------------------------------------------
\multicolumn{2}{l}{\textbf{Neural Tangent Kernel and operators}}\\[0.3em]
$\phi(x),\ \Phi_t(x)$ & Tangent feature map $\nabla_\theta f(x;\theta_0)$ (and its time-$t$ version). \\[0.4em]
$\Theta_0(x,x')$ & Neural Tangent Kernel at initialization, $\nabla_\theta f(x;\theta_0)^\top\nabla_\theta f(x';\theta_0)$. \\[0.4em]
$\Theta_t,\ \Theta_t^{(m)}$ & NTK at training time $t$ (the empirical, width-$m$ kernel). \\[0.4em]
$\Theta_0^{(m)}$ & Finite-width empirical NTK at initialization. \\[0.4em]
$\bar\Theta,\ \Theta$ & Limiting (infinite-width) deterministic NTK; also written as the Gram matrix $\bar\Theta=(\bar\Theta(x_i,x_j))_{i,j}$. \\[0.4em]
$\Theta(\delta)$ & Limiting NTK on $S^1$ as a function of the angular distance $\delta=d(\varphi,\psi)\in[0,\pi]$. \\[0.4em]
$\Theta^{(\mathrm{nb})}$ & Limiting NTK with the trainable hidden-bias contribution removed. \\[0.4em]
$\Theta^{(r)},\ \psi_r$ & One-neuron kernel contribution and one-neuron tangent feature of neuron $r$. \\[0.4em]
$J_0$ & Jacobian of network outputs w.r.t.\ parameters at initialization; $\bar\Theta=J_0 J_0^\ast$. \\[0.4em]
$T_{\bar\Theta},\ T,\ T_\infty$ & Integral operator on $L^2(\mu)$ induced by the limiting kernel; $T_\infty=\lim_{m\to\infty}T_0^{(m)}$. \\[0.4em]
$T_t^{(m)},\ T_0^{(m)}$ & Empirical tangent integral operator at time $t$ and at initialization (frozen operator). \\[0.4em]
$T^{(r)}$ & One-neuron integral operator; $T_0^{(m)}=\tfrac1m\sum_{r=1}^m T^{(r)}$. \\[0.4em]
$T_p$ & Density-weighted integral operator for a non-uniform input density $p(\theta)$. \\[0.4em]
$\kappa$ & The $2\pi$-periodic kernel profile with $\bar\Theta(\theta,\theta')=\kappa(\theta-\theta')$; also a uniform kernel bound and (in the conditioning discussion) a matrix condition number. \\[0.4em]
$r_t,\ r$ & Residual $f_t-y$ (function or vector). \\[0.6em]

% ----------------------------------------------------------------
\multicolumn{2}{l}{\textbf{Fourier basis, spectrum, and subspaces}}\\[0.3em]
$\theta,\ \varphi,\ \gamma$ & Angular coordinate on $S^1$; $\gamma$ is used as the angle variable in target functions. \\[0.4em]
$\phi_q(\theta)=e^{iq\theta}$ & Complex Fourier mode of frequency $q\in\mathbb{Z}$. \\[0.4em]
$\phi_0,\ \phi_{k,c},\ \phi_{k,s}$ & Real Fourier basis of $L^2(\mu)$: $\phi_0=1$, $\phi_{k,c}=\sqrt{2}\cos(k\varphi)$, $\phi_{k,s}=\sqrt{2}\sin(k\varphi)$. \\[0.4em]
$\psi_j$ & Generic orthonormal eigenfunction of an integral operator. \\[0.4em]
$k,\ q$ & Fourier frequency index; $k(p)$ is the frequency of basis function $\phi_p$. \\[0.4em]
$K$ & Frequency cutoff defining the retained low-frequency block. \\[0.4em]
$\mathcal{H}_K$ & Truncated Fourier subspace $\operatorname{span}\{\phi_0,\phi_{k,c},\phi_{k,s}:1\le k\le K\}$, of dimension $d=2K+1$. \\[0.4em]
$\mathcal{F}_k$ & Frequency-$k$ subspace: $\mathcal{F}_0=\operatorname{span}\{\phi_0\}$, $\mathcal{F}_k=\operatorname{span}\{\phi_{k,c},\phi_{k,s}\}$ for $k\ge1$. \\[0.4em]
$d_k$ & Dimension of $\mathcal{F}_k$: $d_0=1$, $d_k=2$ for $k\ge1$. \\[0.4em]
$\lambda_q,\ \lambda_k$ & Fourier eigenvalue of $T$ at frequency $q$ (resp.\ $k$); $\lambda_p:=\lambda_{k(p)}$. \\[0.4em]
$\lambda_k^{(\mathrm{nb})}$ & Fourier eigenvalue of the no-bias kernel $\Theta^{(\mathrm{nb})}$. \\[0.4em]
$\Lambda$ & Diagonal matrix of eigenvalues $\operatorname{diag}(\lambda_p)$; in the linear model, $A=Q\Lambda Q^\top$. \\[0.4em]
$\Lambda_K$ & Continuum reference block $P_K T_\infty P_K=\operatorname{diag}(\lambda_0,\lambda_1,\lambda_1,\dots,\lambda_K,\lambda_K)$. \\[0.4em]
$\alpha_p(t)$ & Amplitude of the residual along Fourier mode $\phi_p$; decays as $e^{-\lambda_p t}\alpha_p(0)$. \\[0.4em]
$a_j(t)$ & Coefficient of the residual along eigenfunction $\psi_j$ in the continuum operator analysis. \\[0.6em]

% ----------------------------------------------------------------
\multicolumn{2}{l}{\textbf{Finite-sample objects}}\\[0.3em]
$\varphi_1,\dots,\varphi_n$ & I.i.d.\ angles sampled uniformly on $[0,2\pi)$. \\[0.4em]
$\Phi$ & Sampled Fourier feature matrix, $\Phi\in\mathbb{R}^{n\times d}$, $\Phi_{i,p}=\phi_p(\varphi_i)$; $\Phi_p$ its $p$-th column. \\[0.4em]
$A$ & Normalized sampled kernel matrix, $A_{ij}=\tfrac1n\Theta(\varphi_i-\varphi_j)$. \\[0.4em]
$G$ & Sampled Fourier Gram matrix, $G=\tfrac1n\Phi^\top\Phi$. \\[0.4em]
$H$ & Compressed operator matrix, $H=\tfrac1n\Phi^\top A\Phi$; $\mathbb{E}[H]=\Lambda^{(n)}$. \\[0.4em]
$\lambda_p^{(n)}$ & Finite-sample corrected Fourier eigenvalue, $(1-\tfrac1n)\lambda_p+\tfrac{\Theta(0)}{n}$. \\[0.4em]
$\Lambda^{(n)}$ & Diagonal matrix $\operatorname{diag}(\lambda_p^{(n)})$ of corrected eigenvalues. \\[0.4em]
$C,\ C_n^{(K)}$ & Coefficient-level restricted operator $G^{-1}H$ on the retained block; compared against $\Lambda^{(n)}$. \\[0.4em]
$P_{\mathrm{th}},\ P_{\mathrm{emp}}$ & Theory and empirical Fourier-block preconditioners. \\[0.4em]
$B_{\mathrm{base}},\ B_{\mathrm{th}},\ B_{\mathrm{emp}}$ & Baseline, theory-preconditioned, and empirically preconditioned sampled operators. \\[0.4em]
$\tau$ & Ridge-regularization parameter for ill-conditioned inverses. \\[0.4em]
$u$ & Unit roundoff (machine precision); the floor $\kappa u$ sets the achievable relative accuracy. \\[0.4em]
$\widehat{r}(t)$ & Sampled residual vector. \\[0.4em]
$\mathbf{1}$ & All-ones vector in $\mathbb{R}^d$. \\[0.6em]

% ----------------------------------------------------------------
\multicolumn{2}{l}{\textbf{Finite-width objects and diagnostics}}\\[0.3em]
$B_m^{(K)}$ & Frozen finite-width tangent block on $\mathcal{H}_K$, $P_K T_0^{(m)} P_K$; $\mathbb{E}[B_m^{(K)}]=\Lambda_K$. \\[0.4em]
$C_t^{(K)}$ & Evolving tangent block on $\mathcal{H}_K$, $P_K T_t^{(m)} P_K$; $C_0^{(K)}=B_m^{(K)}$. \\[0.4em]
$L_t^{(K)}$ & High-to-low frequency coupling (leakage) operator, $P_K T_t^{(m)}(I-P_K)$. \\[0.4em]
$a_t,\ b_t$ & Low- and high-frequency residual components, $a_t=P_K r_t$, $b_t=(I-P_K)r_t$. \\[0.4em]
$\xi_{r,pq}$ & One-neuron matrix entry $\langle\phi_p, T^{(r)}\phi_q\rangle$. \\[0.4em]
$P_K$ & Orthogonal projection $L^2(\mu)\to\mathcal{H}_K$. \\[0.4em]
$P_k,\ U_k$ & Orthogonal projector onto the grid representation of $\mathcal{F}_k$, $P_k=U_kU_k^\top$, with $U_k$ an orthonormal basis. \\[0.4em]
$\widehat{\mathcal{F}}_k(t),\ \widehat P_k(t),\ \widehat U_k(t)$ & Matched empirical eigenspace of the tangent operator, its projector and orthonormal basis. \\[0.4em]
$v_j(t)$ & Orthonormal eigenvectors of the grid-discretized tangent operator. \\[0.4em]
$N,\ \theta_j$ & Number of uniform grid points and grid angles $\theta_j=2\pi j/N$, $j=0,\dots,N-1$. \\[0.4em]
$\varepsilon_{k,F}^{(m)}(t)$ & Projector error $\|\widehat P_k(t)-P_k\|_F$ for frequency $k$ (frozen value at $t=0$ written $\varepsilon_{k,F}^{(m)}$). \\[0.4em]
$\theta_{k,i}(t)$ & $i$-th principal angle between $\widehat{\mathcal{F}}_k(t)$ and $\mathcal{F}_k$; $\cos\theta_{k,i}=s_i(U_k^\top\widehat U_k)$. \\[0.4em]
$s_i(\cdot)$ & $i$-th largest singular value. \\[0.4em]
$s_k^{\mathrm{cos}}(t)$ & Subspace cosine similarity, root-mean-square of the principal-angle cosines for frequency $k$. \\[0.4em]
$\mu_{\mathcal{F}_k}(t)$ & Average spectral strength inside Fourier block $\mathcal{F}_k$, $\tfrac{1}{d_k}\operatorname{tr}(P_k T_t^{(m)} P_k)$. \\[0.4em]
$\mu_{\le K}(t)$ & Cumulative low-frequency spectral strength over the retained block. \\[0.4em]
$\mu_m,\ \mu_t,\ \delta_{\theta_\alpha}$ & Empirical neuron measure, its limit, and the Dirac mass at $\theta_\alpha$ (mean-field discussion). \\[0.6em]

% ----------------------------------------------------------------
\multicolumn{2}{l}{\textbf{Targets and miscellaneous}}\\[0.3em]
$f_{\mathrm{orig}},\ f_{\mathrm{eq}},\ f_{\mathrm{rev}}$ & Original, equal-amplitude, and reversed-amplitude target functions used in the amplitude-variant experiments. \\[0.4em]
$A_i,\ k_i,\ \phi_i$ & Amplitude, frequency, and phase of the $i$-th sinusoid in a synthetic Fourier-regression target. \\[0.4em]
$c$ & Universal constant in the finite-width concentration bound (illustrative value $c=0.2$). \\[0.4em]
$\delta$ & Failure probability in high-probability bounds; also the angular distance $d(\varphi,\psi)$ on $S^1$. \\[0.4em]
$\varepsilon$ & Deviation level in concentration inequalities. \\[0.4em]
$R,\ R_\alpha$ & Rotation matrix on $\mathbb{R}^2$ (rotation by angle $\alpha$). \\

\end{longtable}
